 %%%%%%%%%%%%%%Empieza la presentacion%%%%%%%%%%%%%%%%%%%%%%%%%%%%%%%%%%
 

%%%%%%%%%%%%%%%%%%%%%%%%%%%%%%%%%%%%%%%%%%%%%%%%

\begin{frame}
\frametitle{Potencia de una señal aleatoria}
\vspace{-0.2cm}
\footnotesize
\begin{itemize}
\item La potencia de una señal aleatoria se define$^1$ como el valor esperado de la potencia de una señal:

$P_X=E(\lim_{N \rightarrow \infty } \frac{1}{2N+1}\sum_{n=-N}^{N} x^2(n))=\lim_{N \rightarrow \infty } \frac{1}{2N+1} \sum_{n=-N}^{N} E(x^2(n))$

\visible<2->{

\item En los casos en los que  $E(x^2(n))$ es una señal aleatoria estacionaria se cumple que:

\begin{itemize}
\footnotesize
\item$P_X=E(x^2(n))$ (sumario de 2N+1 términos iguales hace que se cancele el denominador de la expresión).
\item $P_X=\phi_{XX}(0)=\sigma^2_X+\mu_x^2$.
\end{itemize}}


\visible<3->{
\item Por lo tanto, en el caso de las señales aleatorias estacionarias, la potencia de la señal se encuentra relacionada con la autocorrelación en el instante inicial$^2$. 
}

\end{itemize}


\rule{11cm}{0.04ex}
\scriptsize
$^1$Vamos a utilizar las definiciones sobre señales discretas en esta caso. Equivalentemente, se podría utilizar la definición sobre continuas, cambiando el sumatorio por la integral.

\vspace{0.1cm}
\visible<3->{
$^2$Este hecho hace que la autocorrelación se entienda como una potencia generalizada variable con $\tau$.
}

\end{frame}


%%%%%%%%%%%%%%%%%%%%%%%%%%%%%%%%%%%%%%%%%%%%%%%%

\begin{frame}
\frametitle{Transformada de Fourier de señales aleatorias}
\vspace{-0.2cm}
\footnotesize
\begin{itemize}
\footnotesize

\item Las transformadas de Fourier representativas (es decir, que mejor caracterizan) a las señales aleatorias no son la transformada de una muestra/secuencia de la señal, si no  las de las señales de autocorrelacion y correlación cruzada.
\visible<2->{
\item \alert{La densidad espectral de potencia}$^1$  es la transformada de Fourier de la señal de autocorrelación:

\begin{itemize}

\footnotesize

\item Continuo: $\Phi_{XX}(jw)=\int_{-\infty}^{\infty} \phi_{XX}(\tau)e^{-jw\tau}d\tau$

\vspace{0.1cm}

\item Discreto: $\Phi_{XX}(e^{jw})=\sum_{l=-\infty}^{\infty} \phi_{XX}(l)e^{-jwl}$

\vspace{0.1cm}

\item El espectro de potencia permite clasificar las señales aleatorias: ruido blanco, señal aleatoria de banda ancha, señal aleatoria de banda estrecha. 
\end{itemize}

}

\visible<3->{
\item \alert{La densidad espectral cruzada} es la transformada de Fourier de la señal de correlación cruzada:

\begin{itemize}

\footnotesize

\item Continuo: $\Phi_{XY}(jw)=\int_{-\infty}^{\infty} \phi_{XY}(\tau)e^{-jw\tau}d\tau$

\vspace{0.1cm}

\item Discreto: $\Phi_{XY}(e^{jw})=\sum_{l=-\infty}^{\infty} \phi_{XY}(l)e^{-jwl}$

\vspace{0.1cm}

\end{itemize}


}


\end{itemize}


\rule{11cm}{0.04ex}
\scriptsize
También se denomina espectro de densidad de potencia o simplemente espectro de potencia.

\end{frame}


%%%%%%%%%%%%%%%%%%%%%%%%%%%%%%%%%%%%%%%%%%%%%%%%

\begin{frame}
\frametitle{Clasificación de las señales aleatorias}
\vspace{-0.2cm}
\footnotesize
\begin{itemize}
\footnotesize
\item Señal de ruido blanco: aquella cuyo espectro en potencia toma valor constante
\begin{itemize}
\footnotesize
\scriptsize
\item Continuo: $\Phi_{XX}(jw)=a$   $\rightarrow \phi_{XX}(\tau)=a\delta(\tau)$.

\vspace{0.1cm}

 No es físicamente realizable, ya que tendría que ser una señal de potencia infinita (ya que $P_X=\phi_{XX}(\tau=0)=\infty$). 
 
 \vspace{0.1cm}
 
 \item Discreto: $\Phi_{XX}(e^{jw})=a $  $\rightarrow \phi_{XX}(l)=a \delta(l)$.
 
 
\end{itemize}

\item Señal con bias constante: aquella cuya autocorrelación toma un valor constante:
\begin{itemize}
\footnotesize
\scriptsize
\item Continuo: $\Phi_{XX}(jw)=\mu^2_X \delta(w)$   $\leftarrow \phi_{XX}(\tau)=\mu^2_X$.

\vspace{0.1cm}

 \item Discreto: $\Phi_{XX}(e^{jw})=\mu^2_X \delta(w)$   $\leftarrow \phi_{XX}(l)=\mu^2_X$.

 
 \vspace{0.1cm}
 
 
\end{itemize}


 \item Una combinación de las dos anteriores
 
 \item Señal aleatoria con espectro en potencia paso baja: es la que su espectro en potencia tiene la forma de un filtro paso bajo ideal.
 
 \item Señal aleatoria con espectro en potencia paso banda: es la que su espectro en potencia tiene la forma de un filtro paso banda ideal.

\item Señal aleatoria de banda estrecha: aquella cuyo espectro en frecuencia tiene una banda en frecuencias estrecha

\item Señal aleatoria de banda ancha: aquella cuyo espectro en frecuencia tiene una banda en frecuencia ancha.


\end{itemize}
\end{frame}


%%%%%%%%%%%%%%%%%%%%%%%%%%%%%%%%%%%%%%%%%%%%%%%%

\begin{frame}
\frametitle{Procesamiento de señales aleatorias I}
\vspace{-0.2cm}
\footnotesize
\begin{itemize}
\item Las señales aleatorias, como señales que son, se pueden procesar con todos los sistemas lineales y filtros ideales con los que hemos trabajado a lo largo de la asignatura.
\item Pero debido a su aleatoreidad, para caracterizar su comportamiento correctamente necesitamos estudiar que le sucede a la función de probabilidad o a los estadísticos de la señal aleatoria filtrada. 
\item Esto se pude hacer determinando que le sucede a la media, al valor cuadrático medio, autocorrelación, etc ... de la señal de salida del filtro $H(s)$ o $H(z)$ a partir de los valores de los estadísticos correspondientes de la señal de entrada.
\item El análisis se puede realizar tanto en continuo como en discreto. En nuestro caso, realizaremos en discreto el análisis restante del tema.
\item Y en todo momento vamos a suponer que los sistemas lineales/filtros son deterministas, y que por lo tanto, la aleatoreidad del proceso se encuentra únicamente asociado a la aleatoreidad de la señal. 
\item Además, también supondremos que las señales aleatorias son estacionarias (en el sentido amplio).


\end{itemize}
\end{frame}


%%%%%%%%%%%%%%%%%%%%%%%%%%%%%%%%%%%%%%%%%%%%%%%%%


\begin{frame}
\frametitle{Procesamiento de señales aleatorias II}
\vspace{-0.2cm}
\footnotesize
\begin{itemize}
\item La respuesta de un sistema determinista a cualquier entrada se puede calcular a través de la función de convolución: $y(n)=(h*x)(n)=\sum_{k=-\infty}^{\infty} h(k)x(n-k)$ 

\visible<2->{
\item El valor medio de la señal de salida:
\alert{$\mu_Y=H(e^{j0})\mu_X$}

{\scriptsize

\color{gray}
Se calcula aprovechando las propiedades de la linealidad del sistema

$\mu_Y=E(y(n))=E(\sum_{k=-\infty}^{\infty} h(k)x(n-k))=\sum_{k=-\infty}^{\infty} h(k)E(x(n-k))=$

\hspace{0.8cm}$=\mu_X \sum_{k=-\infty}^{\infty} h(k)=H(e^{j0})\mu_X$
\color{black}}}

\visible<3->{

\item La señal de autocorrelación (productos aparecen en el dominio de la frecuencia): \alert{$\Phi_{YY}(e^{jw})=|H(e^{jw})H(e^{-jw})|\Phi_{XX}(e^{jw})$}

{\scriptsize
\color{gray}

$\phi_{YY}(l)\!=\!E(y(n\!+\!l)y(n))\!=\!E\!\!\left (\left ( \sum_{i=-\infty}^{\infty} h(i)x(n\!+\!l\!-\!i) \right )\!\! \left ( \sum_{k=-\infty}^{\infty} h(k)x(n\!-\!k) \right )\right )=$

\hspace{0.9cm}$=\sum_{i=-\infty}^{\infty} h(i) \sum_{k=-\infty}^{\infty} h(k) E(x(n\!+\!l\!-\!i)x(n\!-\!k))=^1$

\vspace{0.1cm}

\hspace{0.9cm}$=\sum_{i=-\infty}^{\infty} h(i) \sum_{k=-\infty}^{\infty} h(k) \psi_{XX}(l+k-i)=^2$

\vspace{0.1cm}

\hspace{0.9cm}$=\sum_{m=-\infty}^{\infty} \psi_{XX}(l-m) \sum_{k=-\infty}^{\infty} h(k) h(m+k) =(\psi_{XX}*(h(k)*h(-k))(l)$

\vspace{0.1cm}

En el dominio de la transformada las operaciones son más sencillas: dos convoluciones sucesivas se convierten en dos productos.



$\phi_{YY}(e^{jw})=|H(e^{jw})H(e^{-jw})|\Phi_{XX}(e^{jw})$

}}

\end{itemize}

\vspace{-0.2cm}

\rule{11cm}{0.04ex}

\visible<3->{
\color{gray}
\scriptsize
$^1$Como son señales estacionarias el valor esperado solo depende de la diferencia de los indices: $E(x(n\!+\!l\!-\!i)x(n\!-\!k))=\psi_{XX}(l+k-i)$

$^2$Cambio de variable $m=i-k$
}


\end{frame}


\begin{frame}
\frametitle{Procesamiento de señales aleatorias III}
\vspace{-0.2cm}
\footnotesize
\begin{itemize}
\item La señal de autocorrelacion: \alert{$\Phi_{YY}(e^{jw})=|H(e^{jw})H(e^{-jw})|\Phi_{XX}(e^{jw})$}
\begin{itemize}
\footnotesize
\item Utilidad: La expresión es especialmente útil para obtener el espectro en frecuencia de un sistema lineal con una función de transferencia desconocida:

\visible<2->{
\begin{itemize}
\footnotesize
\item  Para eso basta con excitar el sistema con una señal aleatoria estacionaria $x(n)$ de espectro en potencia conocido y observar el espectro en potencia de la señal de salida.
 
\item El cociente de ambos espectros nos da el cuadrado del módulo del espectro del sistema lineal.  
\end{itemize}}
\end{itemize}

\visible<3->{
\item La señal de correlación cruzada: \alert{$\Phi_{XY}(e^{jw})=H(e^{jw})\Phi_{XX}(e^{jw})$}

\begin{itemize}
\footnotesize
\item La demostración es análoga a la anterior, calculando la correlación entre la señal de salida (que pasa una vez por el sistema) y la señal de entrada (que no pasa por el filtro). 
\item En este caso, la expresión nos relaciona el espectro de potencia de la señal de entrada con la densidad espectral cruzada de la entrada y la salida. 
\item Utilidad similar a la del caso anterior, pero ahora podemos obtener la respuesta en frecuencia completa.
\end{itemize}}


\end{itemize}
\end{frame}
